TorchDCM implements likelihood evaluation as vectorized PyTorch tensor computation.
Each model maps a data object and parameter vector to utilities, probabilities, and log-likelihood contributions.
Availability masks are applied before normalization, so unavailable alternatives receive zero probability while the feasible set is normalized correctly.
For panel data, row-level probabilities are grouped by individual identifiers and multiplied through the log domain to form individual-level likelihood contributions.

The package uses automatic differentiation for gradients and Hessian-based inference.
Parameter objects define initial values, fixed/free status, and transformations when constraints are needed.
For unconstrained models, optimization proceeds over free parameters directly.
For constrained parameters such as nest dissimilarities, the implementation optimizes an unconstrained internal representation and maps it to the feasible interval during likelihood evaluation.
After estimation, the package computes covariance matrices from observed information where available and exposes standard errors and t-statistics through the result object.

Simulation-based models use explicit draw tensors.
Mixed logit and WTP-space mixed logit average probabilities over deterministic or seeded simulation draws, and panel likelihoods share draws consistently across observations belonging to the same individual.
This design is essential for cross-package validation: when TorchDCM, Biogeme, and Apollo receive the same parameters and the same draw matrix, probability and likelihood differences should isolate implementation errors rather than random simulation noise.
The current benchmark suite uses this design both for mixed-logit likelihood-kernel replay and for an aligned Swissmetro mixed-logit full-estimation comparison against Biogeme.

The validation layer is implemented as executable benchmark scripts rather than static tables.
Each benchmark case records data source, model specification, starting values, availability rules, scaling choices, covariance definition, backend availability, runtime split, and numerical differences.
Generated JSON and Markdown reports are treated as experiment artifacts.
This mirrors the software-paper pattern in which the package is evaluated as a reproducible computational system, not only as a collection of model classes.
